% ID: SEL-D0074-OA-E02
% Tanda S074.
\textbf{a)} El punto de tangencia es $(-1,-3)$. La derivada $f'(x)=-3/x^2$ da la pendiente $f'(-1)=-3$. La ecuación punto-pendiente es $y+3=-3(x+1)$, es decir, $\boxed{y=-3x-6}$.

\textbf{b)} Como el punto pertenece a la gráfica, $g(1)=a+b=2$. La función es derivable cerca de $1$, por lo que en un extremo relativo se cumple $g'(1)=0$. Dado que $g'(x)=a-b/x^2$, resulta $a-b=0$. El sistema da $\boxed{a=b=1}$. Para comprobar que se trata de un extremo, observamos que $g''(1)=2>0$. Por tanto, $(1,2)$ es un mínimo relativo.
